% Options for packages loaded elsewhere
\PassOptionsToPackage{unicode}{hyperref}
\PassOptionsToPackage{hyphens}{url}
\documentclass[
]{article}
\usepackage{xcolor}
\usepackage[margin=1in]{geometry}
\usepackage{amsmath,amssymb}
\setcounter{secnumdepth}{-\maxdimen} % remove section numbering
\usepackage{iftex}
\ifPDFTeX
  \usepackage[T1]{fontenc}
  \usepackage[utf8]{inputenc}
  \usepackage{textcomp} % provide euro and other symbols
\else % if luatex or xetex
  \usepackage{unicode-math} % this also loads fontspec
  \defaultfontfeatures{Scale=MatchLowercase}
  \defaultfontfeatures[\rmfamily]{Ligatures=TeX,Scale=1}
\fi
\usepackage{lmodern}
\ifPDFTeX\else
  % xetex/luatex font selection
\fi
% Use upquote if available, for straight quotes in verbatim environments
\IfFileExists{upquote.sty}{\usepackage{upquote}}{}
\IfFileExists{microtype.sty}{% use microtype if available
  \usepackage[]{microtype}
  \UseMicrotypeSet[protrusion]{basicmath} % disable protrusion for tt fonts
}{}
\makeatletter
\@ifundefined{KOMAClassName}{% if non-KOMA class
  \IfFileExists{parskip.sty}{%
    \usepackage{parskip}
  }{% else
    \setlength{\parindent}{0pt}
    \setlength{\parskip}{6pt plus 2pt minus 1pt}}
}{% if KOMA class
  \KOMAoptions{parskip=half}}
\makeatother
\usepackage{graphicx}
\makeatletter
\newsavebox\pandoc@box
\newcommand*\pandocbounded[1]{% scales image to fit in text height/width
  \sbox\pandoc@box{#1}%
  \Gscale@div\@tempa{\textheight}{\dimexpr\ht\pandoc@box+\dp\pandoc@box\relax}%
  \Gscale@div\@tempb{\linewidth}{\wd\pandoc@box}%
  \ifdim\@tempb\p@<\@tempa\p@\let\@tempa\@tempb\fi% select the smaller of both
  \ifdim\@tempa\p@<\p@\scalebox{\@tempa}{\usebox\pandoc@box}%
  \else\usebox{\pandoc@box}%
  \fi%
}
% Set default figure placement to htbp
\def\fps@figure{htbp}
\makeatother
\setlength{\emergencystretch}{3em} % prevent overfull lines
\providecommand{\tightlist}{%
  \setlength{\itemsep}{0pt}\setlength{\parskip}{0pt}}
\usepackage{bookmark}
\IfFileExists{xurl.sty}{\usepackage{xurl}}{} % add URL line breaks if available
\urlstyle{same}
\hypersetup{
  pdftitle={Association of socio-demographic factors and pain presentation in a cohort of chronic pain patients},
  hidelinks,
  pdfcreator={LaTeX via pandoc}}

\title{Association of socio-demographic factors and pain presentation in
a cohort of chronic pain patients}
\author{}
\date{\vspace{-2.5em}}

\begin{document}
\maketitle

Date:

\subsubsection{Introduction}\label{introduction}

Chronic pain is a complex and multifaceted condition that manifests in
highly heterogeneous ways across individuals. Much variability is
observed in the intensity of pain, the response to conservative
treatment, and scores on other related patient-reported outcome measures
(PROMs) such as anxiety and depression. The aim of this study was to
estimate associations between socio-demographic characteristics and pain
intensity in patients with chronic pain at the intake in a specialized
pain clinic. The secondary aim was to estimate correlations between
PROMs in patients with chronic pain.

\subsubsection{Methods}\label{methods}

For this study, we used a subset of patients from the chronic pain
registry of the Maastricht University Medical Centre department of pain
management who had been recruited between January 2022 and December
2024. This registry has been described in great detail elsewhere.
Patients were considered eligible if they had pain for at least 3 months
and conservative treatment (e.g.~physical therapy) failed. Patients
received questionnaires at enrollment, consisting of several PROMs, such
as the RAND-36 questionnaire and the hospital anxiety and depression
scale (HADS). Baseline data were summarized as mean and standard
deviation (SD) or as count and percentage. Multivariable linear
regression analysis was used to estimate associations between
socio-demographic factors and pain intensity. Pearson's correlation
coefficient was computed to estimate the correlation between pain
intensity, depressive complaints, and anxiety.

\subsubsection{Results}\label{results}

\paragraph{Baseline characteristics}\label{baseline-characteristics}

Between January 2022 and December 2024, we included chronic pain
patients, of whom (\%) were female, with a mean age (SD) of (). Table 1
shows baseline characteristics of the cohort of chronic pain patients. `

\paragraph{Associations between socia-demographic factors and pain
intensity}\label{associations-between-socia-demographic-factors-and-pain-intensity}

The mean (SD) pain intensity measured using the RAND questionnaire pain
domain was (). Linear regression analysis did not reveal any
statistically significant associations between socio-demographic factors
and pain intensity (table 2). However, both female sex (regression
coefficient = ) and having a child (regression coefficient = ) were
associated with a clinically meaningful higher mean pain intensity,
albeit not significant.

\paragraph{Correlations between outcome
measures}\label{correlations-between-outcome-measures}

Figure 1 shows the correlations between different outcome measures.
Strong negative correlations were seen between the HADS depression and
anxiety scales on the one hand, and the RAND questionnaire mental health
dimension, and moderate positive correlations between the two HADS
subscales. The RAND pain dimension was not strongly correlated to any of
the other patient-reported outcomes.

\subsubsection{Conclusion}\label{conclusion}

We did not find strong significant associations between
socio-demographic factors and pain intensity in a large cohort of
chronic pain patients, albeit some were potentially clinically relevant.
Pain intensity was not strongly correlated to other patient outcomes.

\subsection{References}\label{references}

\end{document}
